% from ega0-1-fr.tex
\oldpage[0\textsubscript{I}]{48}

\begin{env}[5.3.8]
\label{0.5.3.8-fr}
Supposons $\mathcal{O}_X$ cohérent~; soient $\mathcal{F}$ un
$\mathcal{O}_X$-Module cohérent, $x$ un point de $X$, $M$ un sous-module de
type fini de $\mathcal{F}_x$~; il existe alors un voisinage ouvert $U$ de $x$
et un sous-$(\mathcal{O}_X|U)$-Module cohérent $\mathcal{G}$ de
$\mathcal{F}|U$ tel que $\mathcal{G}_x=M$ (T, 4.1, lemme 1).

Ce résultat, joint aux propriétés des sous-$\mathcal{O}_X$-Modules d'un
$\mathcal{O}_X$-Module cohérent, impose des conditions nécessaires aux anneaux
$\mathcal{O}_x$ pour que $\mathcal{O}_X$ soit cohérent. Par exemple (5.3.4),
l'intersection de deux idéaux de type fini de $\mathcal{O}_x$ doit encore être
un idéal de type fini.
\end{env}

\begin{env}[5.3.9]
\label{0.5.3.9-fr}
Supposons $\mathcal{O}_X$ cohérent, et soit $M$ un $\mathcal{O}_x$-module
admettant une présentation finie, donc isomorphe à un conoyau d'un
homomorphisme $\varphi:\mathcal{O}_x^p\to\mathcal{O}_x^q$~; alors il existe un
voisinage ouvert $U$ de $X$ et un $(\mathcal{O}_X|U)$-Module cohérent
$\mathcal{F}$ tel que $\mathcal{F}_x$ soit isomorphe à $M$. En effet, en vertu
de (5.2.6), il existe une section $u$ de
$\mathcal{H}\!om_{\mathcal{O}_X}(\mathcal{O}_X^p,\mathcal{O}_X^q)$ au-dessus
d'un voisinage ouvert $U$ de $x$ telle que $u_x=\varphi$~; le conoyau
$\mathcal{F}$ de l'homomorphisme
$u:\mathcal{O}_X^p|U\to\mathcal{O}_X^q|U$ répond à la question (5.3.4).
\end{env}

\begin{env}[5.3.10]
\label{0.5.3.10-fr}
Supposons $\mathcal{O}_X$ cohérent, et soit $\mathcal{J}$ un faisceau cohérent
d'idéaux dans $\mathcal{O}_X$. Pour qu'un
$(\mathcal{O}_X/\mathcal{J})$-Module $\mathcal{F}$ soit cohérent, il faut et il
suffit qu'il soit cohérent en tant que $\mathcal{O}_X$-Module. En particulier
$\mathcal{O}_X/\mathcal{J}$ est un faisceau cohérent d'anneaux.
\end{env}

\begin{env}[5.3.11]
\label{0.5.3.11-fr}
Soit $f:X\to Y$ un morphisme d'espaces annelés, et supposons
$\mathcal{O}_X$ cohérent~; alors, pour tout $\mathcal{O}_Y$-Module cohérent
$\mathcal{G}$, $f^*(\mathcal{G})$ est un $\mathcal{O}_X$-Module cohérent. En
effet, avec les notations de (5.2.4), on peut supposer que $\mathcal{G}|V$ est
conoyau d'un homomorphisme
$v:\mathcal{O}_Y^q|V\to\mathcal{O}_Y^p|V$~; comme $f_U^*$ est exact à droite,
$f^*(\mathcal{G})|U=f_U^*(\mathcal{G}|V)$ est conoyau de l'homomorphisme
$f_U^*(v):\mathcal{O}_X^q|U\to\mathcal{O}_X^p|U$, d'où notre assertion.
\end{env}

\begin{env}[5.3.12]
\label{0.5.3.12-fr}
Soient $Y$ une partie fermée de $X$, $j:Y\to X$ l'injection canonique,
$\mathcal{O}_Y$ un faisceau d'anneaux sur $Y$, et posons
$\mathcal{O}_X=j_*(\mathcal{O}_Y)$. Pour qu'un $\mathcal{O}_Y$-Module
$\mathcal{G}$ soit de type fini (resp. quasi-cohérent, cohérent), il faut et il
suffit que $j_*(\mathcal{G})$ soit un $\mathcal{O}_X$-Module de type fini
(resp. quasi-cohérent, cohérent).
\end{env}

\subsection{Faisceaux localement libres}
\label{subsection:0.5.4-fr}

\begin{env}[5.4.1]
\label{0.5.4.1-fr}
Soit $X$ un espace annelé. On dit qu'un $\mathcal{O}_X$-Module $\mathcal{F}$
est \emph{localement libre} si, pour tout $x\in X$, il existe un voisinage
ouvert $U$ de $x$ tel que $\mathcal{F}|U$ soit isomorphe à un
$(\mathcal{O}_X|U)$-Module de la forme $\mathcal{O}_X^{(I)}|U$, où $I$ peut
dépendre de $U$. Si pour tout $U$, $I$ est fini, on dit que $\mathcal{F}$ est
\emph{de rang fini}~; si pour tout $U$, $I$ a le même nombre fini d'éléments
$n$, on dit que $\mathcal{F}$ est \emph{de rang $n$}. Un
$\mathcal{O}_X$-Module localement libre de rang 1 est encore appelé
\emph{inversible} (cf. (5.4.3)). Si $\mathcal{F}$ est un
$\mathcal{O}_X$-Module localement libre de rang fini, pour tout $x\in X$,
$\mathcal{F}_x$ est un $\mathcal{O}_x$-module libre de rang fini $n(x)$, et il
existe un voisinage $U$ de $x$ tel que $\mathcal{F}|U$ soit de rang $n(x)$~;
si $X$ est connexe, $n(x)$ est donc \emph{constant}.

Il est clair que tout faisceau localement libre est quasi-cohérent, et si
$\mathcal{O}_X$ est un faisceau cohérent d'anneaux, tout
$\mathcal{O}_X$-Module localement libre de rang fini est cohérent.

Si $\mathcal{L}$ est localement libre,
$\mathcal{L}\otimes_{\mathcal{O}_X}\mathcal{F}$ est un foncteur \emph{exact}
en $\mathcal{F}$ dans la catégorie des $\mathcal{O}_X$-Modules.

Nous aurons surtout à considérer des $\mathcal{O}_X$-Modules localement libres
de rang fini,

